%! Representing a Quantitative Variable with Graphs — Unit 1, Lesson 1.3 — Experience & Formalize
%! (Activity · QuickNotes · Application · Check Your Understanding) (Answer Key)
% EFFL component, Math Medic sizing: 12pt body, OPEN white space after every question.
% Page budget: Activity <=2pp · QuickNotes 1/2pp · Application 1/2-1pp · CYU 1-2pp.
%
% SPOILER RULE: the Activity never says histogram, dotplot, stemplot, discrete, continuous, or
% distribution. The displays are called "the plot" and "the two pictures". `Categorical' and
% `quantitative' WERE formalized in Lesson 1.1 and are prior knowledge here.
%
% NO SKETCH QUESTIONS (hard constraint): every display is pre-drawn. Students "construct" only
% by filling in the interval table in 1c, never by drawing an axis or plotting a point.
%
% THE DATA (24 quiz finishing times, used in every display in this lesson):
%   11 12 12 14 15 15 16 17 18 18 19 20 | 34 35 36 36 38 39 40 41 42 43 44 46
% Width-5 bins from 10: 4, 7, 1, 0, 1, 5, 5, 1  -> two clear clumps with a gap at 25-30.
% Width-20 bins from 10: 12, 12            -> two equal bars; the gap vanishes completely.
% That collapse is the crux (EK UNC-1.G.1: altering interval widths changes the appearance).
% GENERATED FROM experience/main.tex — the blank with answers filled in, so every \answerspace
% height and \boxguard is identical to it. NO teachernote here — that lives in the lesson plan.
\documentclass[12pt]{article}
\usepackage{apstats-article}
\usepackage{apstats-key}   % pulls in -boxes; do NOT also load -boxes
\newcommand{\answerspace}[2]{\par\nopagebreak\noindent\begin{minipage}[t][#1][t]{\linewidth}%
  \color{keyred}\bfseries #2\end{minipage}\par}

\begin{document}

\pageheader{Unit 1, Lesson 1.3}{Experience \& Formalize: Twenty-Four Finishing Times --- Answer Key}
% NAMESTRIP: no \namedateperiod here — mirrors the blank.

\vspace{0.02in}

% ── Part 1: the Activity ─────────────────────────────────────────────────────
\begin{headlinebox}{sky}
\textbf{Twenty-Four Finishing Times.} A teacher gave a 50-minute quiz and wrote down how many
minutes each of her 24 students took to finish. The same 24 numbers appear three times below,
drawn three different ways. Work through the two problems with your group using only what you
already know.
\end{headlinebox}

\vspace{0.04in}

\begin{scenariobox}[1. The Class Record]{navy}
Here is how the teacher first wrote them down. The number to the left of the line is the
\emph{tens} digit; each number to the right is a \emph{ones} digit.

\vspace{5pt}
\begin{center}
\renewcommand{\arraystretch}{1.25}
\begin{tabular}{@{} r | l @{}}
\textbf{1} & 1 \quad 2 \quad 2 \quad 4 \quad 5 \quad 5 \quad 6 \quad 7 \quad 8 \quad 8 \quad 9 \\
\textbf{2} & 0 \\
\textbf{3} & 4 \quad 5 \quad 6 \quad 6 \quad 8 \quad 9 \\
\textbf{4} & 0 \quad 1 \quad 2 \quad 3 \quad 4 \quad 6 \\
\end{tabular}
\end{center}
\vspace{3pt}

\normalsize
\begin{enumerate}[label=\textbf{\alph*.}, leftmargin=1.8em, itemsep=6pt]
  \item The row marked \textbf{3} stands for the times 34, 35, 36, 36, 38, and 39.

        How many students finished in total? \ans{24} \quad
        Fastest time: \ans{11 min} \quad Slowest: \ans{46 min}
  \item Could one of these students have finished in exactly \textbf{15.3} minutes? \quad
        Could the same teacher have recorded \textbf{15.3} as the \emph{number of questions a
        student attempted}? Explain what makes those two different.
        \answerspace{2.0cm}{Yes to 15.3 minutes --- a stopwatch can read any value, and between
        any two times there is always another. No to 15.3 questions: you attempt 15 or 16 and
        nothing between. One is measured, the other counted.}
  \item Group the 24 times into five-minute blocks by counting them off the plot above.

        \vspace{3pt}
        \renewcommand{\arraystretch}{1.3}
        \begin{tabularx}{\linewidth}{@{} l Y Y Y Y Y Y Y Y @{}}
        \rowcolor{sky}
        \textbf{Minutes} & 10--14 & 15--19 & 20--24 & 25--29 & 30--34 & 35--39 & 40--44 & 45--49 \\
        \hline
        \textbf{Students} & \ans{4} & \ans{7} & \ans{1} & \ans{0} &
        \ans{1} & \ans{5} & \ans{5} & \ans{1} \\
        \end{tabularx}
  \item Your table and the plot above hold the same 24 students. Name one thing the plot can
        tell you that the table no longer can --- and one thing the table makes easier to see.
        \answerspace{2.0cm}{The plot gives back every individual time --- you can still read the
        two 36s. The table only says ``5 students between 35 and 39.'' But the table makes the
        overall pattern obvious: two clusters with an empty stretch between them.}
\end{enumerate}
\end{scenariobox}

\boxguard[16]
\begin{scenariobox}[2. Two Pictures of the Same 24 Times]{navy}
Both pictures below were drawn from the \emph{same} 24 finishing times you just worked with.
Nothing was left out of either one.

\vspace{5pt}
\begin{center}
\begin{tikzpicture}[scale=0.86]
  % ── left: five-minute blocks ──
  \foreach \y/\lab in {0/0, 1.4/5, 2.8/10} {
    \draw[linegray, very thin] (0,\y) -- (5.6,\y);
    \draw (-0.08,\y) -- (0.08,\y) node[left, font=\scriptsize, xshift=-2pt] {\lab};}
  \foreach \i/\c in {0/4, 1/7, 2/1, 3/0, 4/1, 5/5, 6/5, 7/1} {
    \fill[navy] ({\i*0.7},0) rectangle ({(\i+1)*0.7},{\c*0.28});
    \draw[white, very thin] ({\i*0.7},0) rectangle ({(\i+1)*0.7},{\c*0.28});}
  \draw[->, thick] (0,0) -- (0,4.0);
  \draw[->, thick] (0,0) -- (6.0,0);
  \foreach \x/\lab in {0/10, 1.4/20, 2.8/30, 4.2/40, 5.6/50} {
    \draw (\x,-0.08) -- (\x,0.08); \node[below, font=\scriptsize] at (\x,-0.12) {\lab};}
  \node[font=\scriptsize] at (2.8,-0.75) {Minutes to finish};
  \node[font=\bfseries\small] at (2.8,4.3) {Blocks of 5 minutes};
  % ── right: twenty-minute blocks ──
  \begin{scope}[xshift=8.2cm]
    \foreach \y/\lab in {0/0, 1.4/5, 2.8/10} {
      \draw[linegray, very thin] (0,\y) -- (5.6,\y);
      \draw (-0.08,\y) -- (0.08,\y) node[left, font=\scriptsize, xshift=-2pt] {\lab};}
    \foreach \i/\c in {0/12, 1/12} {
      \fill[goldacc] ({\i*2.8},0) rectangle ({(\i+1)*2.8},{\c*0.28});
      \draw[white, very thin] ({\i*2.8},0) rectangle ({(\i+1)*2.8},{\c*0.28});}
    \draw[->, thick] (0,0) -- (0,4.0);
    \draw[->, thick] (0,0) -- (6.0,0);
    \foreach \x/\lab in {0/10, 2.8/30, 5.6/50} {
      \draw (\x,-0.08) -- (\x,0.08); \node[below, font=\scriptsize] at (\x,-0.12) {\lab};}
    \node[font=\scriptsize] at (2.8,-0.75) {Minutes to finish};
    \node[font=\bfseries\small] at (2.8,4.3) {Blocks of 20 minutes};
  \end{scope}
\end{tikzpicture}
\end{center}
\vspace{4pt}

\normalsize
\begin{enumerate}[label=\textbf{\alph*.}, leftmargin=1.8em, itemsep=6pt]
  \item Add up the bar heights in each picture. Left: \ans{24} \quad
        Right: \ans{24} \quad Should they match? \ans{Yes --- same 24}
  \item Write one sentence describing what the \textbf{left} picture says about when this class
        finished, and one sentence for the \textbf{right}.
        \answerspace{2.2cm}{Left: the class split in two --- one group finished around 11--20
        minutes, another around 34--46, and almost nobody finished in between.\\[2pt]
        Right: students finished fairly evenly across the whole period, with no pattern.}
  \item \textbf{Both pictures are correct, and they tell different stories.} Say exactly what
        was changed to turn one into the other. Then say which one you would hand the teacher
        who wants to know whether the quiz was the right length --- and why.
        \answerspace{2.6cm}{Only the width of the blocks changed --- from 5 minutes to 20.
        Nothing was added or removed. Hand her the left one: it shows the class splits into fast
        and slow finishers with a real gap, which is exactly what she needs to judge the length.
        The 20-minute blocks average that gap away and make the class look uniform.}
  \item Suppose a third picture used blocks of \textbf{one} minute. Predict what it would look
        like, and say whether it would help this teacher more or less than the left one.
        \answerspace{2.0cm}{Almost every bar would be 0 or 1 tall, since the 24 times barely
        repeat. It would show every value but no shape --- less helpful than the left one. Too
        narrow shows noise, too wide hides the pattern.}
\end{enumerate}
\end{scenariobox}

% ── Part 2: QuickNotes (the debrief fills this) — target: half a page ────────
\boxguard[14]
\begin{tcolorbox}[enhanced, colback=sky, colframe=navy, boxrule=0.8pt, arc=2mm,
    left=3mm, right=3mm, top=1.5mm, bottom=1.5mm, breakable,
    fonttitle=\bfseries\small\color{white}, title={QuickNotes: Displays for a Quantitative Variable},
    attach boxed title to top left={yshift=-2mm, xshift=4mm},
    boxed title style={colback=navy, colframe=navy, arc=1.5mm}]
\small
\begin{itemize}[leftmargin=1.1em, itemsep=4pt]
  \item A \textbf{discrete} variable takes a \ans{countable} number of values (siblings,
        questions attempted). A \textbf{continuous} variable can take \ans{infinitely} many ---
        between any two values there is always another (time, height).
  \item \textbf{Dotplot:} one dot per observation, stacked. \textbf{Stemplot:} each value split
        into a \ans{stem} (leading digits) and a \ans{leaf} (last digit). Both keep
        \ans{every} individual value.
  \item \textbf{Histogram:} the height of each bar is the number of observations falling in
        that \ans{interval}. Individual values are \ans{lost}.
  \item \textbf{Bars touch} in a histogram, because the horizontal axis is a
        \ans{number line}. (In a bar graph the categories were separate, so those bars had gaps.)
  \item \textbf{Changing the interval width changes the \ans{appearance} of the histogram} ---
        without changing a single data value. Too wide hides features; too narrow shows noise.
  \item A \textbf{cumulative graph} shows how many observations are \ans{less than or equal to} a given value.
\end{itemize}
\end{tcolorbox}

% ── Part 3: the Application (worked together, ~6 min) ────────────────────────
\boxguard[14]
\begin{notesbox}{Application: The Bus Stop}
We will do this one together. Here are the waiting times, in minutes, for 18 riders at one bus
stop.

\vspace{3pt}
\begin{center}
\begin{tikzpicture}[scale=0.9]
  \draw[->, thick] (-0.2,0) -- (7.4,0);
  \foreach \x/\lab in {0/0, 1/2, 2/4, 3/6, 4/8, 5/10, 6/12, 7/14} {
    \draw (\x,-0.07) -- (\x,0.07); \node[below, font=\scriptsize] at (\x,-0.1) {\lab};}
  \foreach \x/\n in {0.5/2, 1/3, 1.5/4, 2/3, 2.5/2, 3/1, 5.5/1, 6/2} {
    \foreach \k in {1,...,\n} { \fill[navy] (\x,{0.22*\k}) circle (0.075); }}
  \node[font=\scriptsize] at (3.6,-0.6) {Minutes waited};
\end{tikzpicture}
\end{center}
\vspace{2pt}

\begin{enumerate}[label=\textbf{\alph*.}, leftmargin=1.8em, itemsep=6pt]
  \item How many riders are shown? \ans{18} \quad How many waited more than 8 minutes?
        \ans{3 (at 11 and 12)}
  \item Is waiting time discrete or continuous? \ans{Continuous} \quad Explain how you can tell
        from the variable itself, not from the picture.
        \answerspace{1.6cm}{Between any two waits there is always another --- 3.5 minutes sits
        between 3 and 4. Time is measured, not counted.}
  \item \textbf{What if we redraw it?} If these 18 times were put into blocks of 5 minutes
        instead, what would happen to the two riders out at 11 and 12 minutes --- and would the
        picture still show that they stand apart?
        \answerspace{1.8cm}{They would land together in the 10--15 block, a small bar of height
        3. They would still show as separate, but the empty stretch from 6 to 11 minutes --- the
        thing that makes them look unusual --- would be hidden.}
\end{enumerate}
\end{notesbox}

% ── Part 4: Check Your Understanding — practice, NOT scored ──────────────────
\boxguard[14]
\begin{notesbox}{Check Your Understanding}
Work with your partner, then finish on your own. \emph{These are for practice --- they are not
scored.}
\begin{enumerate}[leftmargin=1.9em, itemsep=6pt]
  \item Label each variable \textbf{discrete} or \textbf{continuous}.

        \vspace{3pt}
        \renewcommand{\arraystretch}{1.3}
        \begin{tabularx}{\linewidth}{@{} Y Y Y Y @{}}
        \textbf{Cars in a lot} & \textbf{Height (cm)} & \textbf{Goals scored} &
        \textbf{Time to run 400 m} \\
        \ans{discrete} & \ans{continuous} & \ans{discrete} & \ans{continuous} \\
        \end{tabularx}
  \item A histogram of 60 test scores uses blocks of 10 points and shows one tall middle bar.
        Redrawn with blocks of 2 points, it shows two separate peaks. Which version is
        \emph{wrong}? Explain.
        \answerspace{1.8cm}{Neither. Both are correct displays of the same 60 scores. Changing
        the block width changed what is visible: the 2-point version reveals two peaks that the
        10-point version averaged away.}
  \item In a histogram the bars touch; in the bar graph you drew last lesson they did not.
        Explain what that difference is telling you about the horizontal axis.
        \answerspace{1.8cm}{A histogram's axis is a number line, so the intervals are adjacent
        and the bars must touch. A bar graph's axis holds separate categories with no order and
        no distance between them, so the gaps say they are not a continuum.}
  \item A stemplot has stem \textbf{7} with leaves \textbf{2 4 4 9}. Write out the four values
        those represent. \ans{72, \quad 74, \quad 74, \quad 79}

        Why can a histogram of the same data not give you those four numbers back?
        \answerspace{1.4cm}{A histogram records only how many values fall in each interval, so
        grouping discards the individual numbers.}
  \item \textbf{Formative check.} A researcher has 200 reaction times measured to the nearest
        thousandth of a second and wants a display that shows the overall shape. Which
        statement is \textbf{correct}? Circle one, then justify it in a sentence.
        \begin{enumerate}[label=\textbf{(\Alph*)}, leftmargin=2.2em, itemsep=2pt]
          \item A stemplot is best, because it keeps every individual value.
          \item \ans{A histogram is reasonable, because with 200 values the shape matters more than
                the individual numbers.}
          \item A bar graph is best, because reaction time is measured numerically.
          \item No display works, because the values are continuous.
          \item A histogram is wrong, because it changes appearance with the interval width.
        \end{enumerate}
        \answerspace{1.4cm}{(B). A stemplot of 200 continuous values would be unreadable, and
        the researcher asked about shape, not individual times.}
\end{enumerate}
\end{notesbox}

\end{document}
